\subsection{Analisis Penegakan Otoritas pada \textit{Server} PRE}
\label{subsec:analisis-penegakan-otoritas-pre}

Pada DecMed, PRE digunakan untuk mengubah data terenkripsi dari pasien ke penerima akses tanpa membuka \textit{plaintext}. PRE tidak menentukan sendiri apakah pengguna berhak mengakses data; keputusan otorisasi ditentukan oleh \textit{policy} PBAC dan token Macaroon. PRE memungkinkan data atau kunci terenkripsi dibuka oleh penerima akses yang sah dan mendukung delegasi kemampuan dekripsi tanpa membagikan kunci privat pasien secara langsung.

\textit{Server} PRE dipilih sebagai titik penegakan karena berada pada jalur kritis sebelum data terenkripsi diproses. Pengguna tidak memperoleh \textit{plaintext} langsung dari penyimpanan \textit{off-chain}, melainkan membutuhkan proses \textit{re-encryption} yang dikendalikan oleh PRE. Pada jalur tersebut, \textit{Server} PRE memverifikasi token dan atribut terkait, menolak permintaan yang tidak sah, lalu melanjutkan \textit{re-encryption} hanya jika otorisasi valid.

\subsubsection{Analisis Verifikasi Aktor dan Operasi Tulis}
\label{subsec:analisis-verifikasi-aktor-operasi-tulis}

Pada DecMed sebelumnya, \textit{role} \texttt{MedicalPersonnel} digunakan untuk merepresentasikan tenaga kesehatan klinis secara umum dan memberi hak mengakses data klinis pasien. \textit{Role} tersebut belum cukup untuk membedakan dokter, perawat, petugas laboratorium, dan apoteker, padahal keempat aktor memiliki kewenangan tulis yang berbeda. Sebagai contoh, dokter dapat menulis diagnosis dan terapi, sedangkan petugas laboratorium hanya dapat menulis hasil laboratorium.

Diperlukan verifikasi tambahan terhadap \textit{subrole}, terutama untuk operasi tulis. \textit{Server} PRE memastikan bahwa aktor pemohon memiliki \textit{subrole} yang sesuai dengan \textit{dataset category} yang akan ditulis. Verifikasi ini menolak kondisi ketika aktor membawa token dengan izin tulis, tetapi secara kewenangan profesional tidak berhak memperbarui segmen data tertentu.

\subsubsection{Analisis Verifikasi Bukti Identitas Aktor pada Rantai Delegasi}
\label{subsec:analisis-verifikasi-bukti-identitas-aktor-aktif}

Pada DecMed sebelumnya, proses otorisasi berfokus pada pemeriksaan klaim yang terdapat pada token. Namun, mekanisme tersebut belum memastikan bahwa pengirim \textit{request} merupakan aktor yang identitasnya tercantum pada token. Kondisi ini menimbulkan kebutuhan verifikasi tambahan, yaitu verifikasi bukti kepemilikan identitas aktor.

Identitas pengguna pada DecMed direpresentasikan melalui \textit{wallet} yang dibuat saat pendaftaran akun. \textit{Wallet address} menjadi identifikator unik yang dapat membedakan setiap pengguna pada sistem. Oleh karena itu, \textit{wallet} digunakan sebagai dasar verifikasi identitas aktor yang menggunakan token.

Pada token hasil delegasi, himpunan aktor yang berwenang dibentuk dari \texttt{root\_subject} dan seluruh \texttt{delegated\_to} pada rantai delegasi. Verifikasi dilakukan dengan meminta pemohon membuktikan kepemilikan \textit{wallet} yang cocok dengan salah satu aktor pada himpunan tersebut. Dengan mekanisme ini, setiap \textit{request} membawa bukti bahwa pengirim \textit{request} menguasai \textit{wallet} dari aktor yang tercatat. Tanda tangan dari aktor di luar rantai ditolak. Pembuktian tersebut ditujukan kepada \textit{Server} PRE sebagai penegak \textit{policy} akses. Hasil verifikasi ini kemudian digunakan oleh \textit{Server} PRE untuk menentukan apakah pemohon berwenang melanjutkan operasi akses terhadap data RME.

\subsubsection{Analisis Revokasi dan Audit Delegasi}
\label{subsec:analisis-revokasi-audit-delegasi}

Delegasi akses menyebabkan token tidak lagi berdiri sendiri, tetapi membentuk relasi antara token induk dan token turunan. Oleh karena itu, mekanisme revokasi pada DecMed perlu diperluas agar pencabutan token induk dapat berdampak pada token turunan yang berasal dari token tersebut. Tanpa pencabutan berantai, aktor yang memperoleh token turunan masih dapat menggunakan akses meskipun akses asalnya telah dicabut oleh pasien.

Revokasi berantai diperlukan untuk mempertahankan kendali pasien sebagai pemilik data. Ketika pasien mencabut akses awal, sistem harus dapat mengenali bahwa token-token turunan yang berada dalam rantai delegasi yang sama tidak lagi memiliki dasar kewenangan yang sah. Dengan demikian, pemeriksaan revokasi tidak cukup hanya dilakukan terhadap token yang dikirim pada saat permintaan akses, tetapi juga perlu mempertimbangkan hubungan token tersebut dengan token asalnya.

Selain revokasi, delegasi juga membutuhkan audit yang mencatat hubungan antaraktor. Audit tidak cukup hanya mencatat bahwa suatu segmen diakses, tetapi juga perlu mencatat sumber kewenangan akses tersebut. Informasi ini diperlukan agar pasien dan sistem dapat menelusuri bagaimana akses terhadap data RME terbentuk selama proses pelayanan.